Phase I included a pre-specified value-prediction protocol on its admitted
populations. Its outcome is a boundary result for the routing literature;
we state its scope precisely, because the Phase-I design (task overlap,
proxy judge, two populations) limits what it can establish.

\paragraph{Open-set value prediction does not beat the best fixed harness.}
A leave-one-harness-out (LOHO) predictor --- logistic heads on task TF-IDF
plus harness code and structural features --- fails to beat the best fixed
harness on both admitted Phase-I populations (DeepSeek arm $-3.4$\,pp,
Qwen arm $+0.9$\,pp, bootstrap CI $[-5.2, +6.9]$; cross-protocol pool
$-0.9$\,pp, CI $[-6.0, +4.3]$). Apparent per-fold repair AUROC of
0.82--0.98 collapses under task-held-out controls: with
database-stratified 50/50 task splits (five seeds, task text + harness
identity features), repair classification is at chance (mean AUROC 0.512,
seed range 0.35--0.60) and the induced policy matches bare ($+0.27$\,pp).
Within this tested configuration, the apparent signal is task
memorization, not harness-value estimation. Unseen-harness representations
add $\leq$0.001 AUROC over task-only features.

\paragraph{The one positive result is narrow.} On the 24-harness
cross-protocol pool only, the same predictor used as an \emph{abstaining
gate} gained $+6.9$\,pp over bare with zero \emph{observed} induced harm
on its interventions. We state its scope: it is measured on the same tasks
the gate was tuned on, on a pool assembled post hoc across protocols;
zero observed harm on these items is not a guarantee of zero harm on
deployment, and single-builder populations showed intervention harm of
2.8\% and 8.7\%. Selective abstention may be easier than full routing;
this observation motivates, but does not establish, that claim.

We conclude: behavioral diversity is necessary for routing value to
exist, and estimating an unseen harness's per-item value from code and
task text alone remains open --- on the predictors, tasks, and splits
tested here.
